\documentclass[UTF8,a4paper,11pt]{ctexart}
\usepackage[margin=22mm,headheight=14pt,footskip=12mm]{geometry}
\usepackage{amsmath,amssymb,array,longtable,booktabs,tikz,needspace,fancyhdr,hyperref}
\usepackage{microtype}
\hypersetup{colorlinks=false,pdfborder={0 0 0},pdfauthor={Good Learning},bookmarks=false}
\newfontfamily\circledfont{Noto Sans Symbols 2}
\newcolumntype{P}[1]{>{\raggedright\arraybackslash}p{#1}}
\setlength{\LTleft}{0pt}\setlength{\LTright}{0pt}
\setlength{\parindent}{0pt}\setlength{\parskip}{.65em}
\linespread{1.18}
\raggedbottom
\setcounter{secnumdepth}{0}
\ctexset{section={format=\Large\bfseries\sffamily,beforeskip=1.35em,afterskip=.65em},
subsection={format=\large\bfseries\sffamily,beforeskip=1em,afterskip=.45em}}
\newcommand{\chaptertitle}[1]{\Needspace{8\baselineskip}\noindent{\LARGE\bfseries\sffamily #1}\par\vspace{4mm}\hrule\vspace{6mm}}
\newcommand{\CompactQuestion}[2]{\par\noindent\textbf{#1.}\ #2\par\vspace{.8em}}
\newcommand{\PracticeQuestion}[3]{\par\Needspace{4\baselineskip}\noindent\begin{minipage}{\linewidth}\textbf{#1.}\ #2\par\noindent\rule{0pt}{#3mm}\par\end{minipage}\par}
\newcommand{\CompactQuestionColon}[2]{\par\noindent\textbf{#1：}\ #2\par\vspace{.8em}}
\newcommand{\PracticeQuestionColon}[3]{\par\Needspace{4\baselineskip}\noindent\begin{minipage}{\linewidth}\textbf{#1：}\ #2\par\noindent\rule{0pt}{#3mm}\par\end{minipage}\par}
\pagestyle{fancy}\fancyhf{}\fancyfoot[C]{\thepage}\renewcommand{\headrulewidth}{0pt}
\begin{document}
\thispagestyle{empty}
\noindent{\sffamily\footnotesize 每日课程}\hfill{\sffamily\footnotesize 学习导图}\par
\vspace{7mm}
\noindent{\sffamily\fontsize{25}{30}\selectfont\bfseries 第 2 天 · 绪论与数学基础（紧凑排版型）}\par
\vspace{4mm}\noindent\rule{\linewidth}{.85pt}\par\vspace{6mm}
\begin{center}
\begin{tikzpicture}[x=1cm,y=1cm,>=stealth,every node/.style={font=\sffamily}]
\path[use as bounding box] (-8.1,-18.7) rectangle (8.1,.4);
\node[font=\sffamily\Large\bfseries] at (0,-.6) {人工智能课程的四项数学工具};
\node[font=\sffamily\small,text=black!55] at (0,-1.25) {2026 年 11 月 16 日 · K01–K02};
\node[anchor=north west,font=\sffamily\fontsize{22}{24}\selectfont\bfseries,text=black!50] at (-7.35,-2.35) {01};
\node[anchor=north west,text width=5.8cm,align=left,font=\sffamily\large\bfseries] at (-6.199999999999999,-2.3200000000000003) {人工智能绪论};
\node[anchor=north west,text width=5.9cm,align=left,font=\sffamily\small] at (-6.199999999999999,-3.0700000000000003) {图灵测试评估文字对话表现\\[.25em]视觉处理图像，语言处理文本};
\draw[black!35,line width=.4pt] (-7.35,-5.050000000000001) -- (-0.2999999999999998,-5.050000000000001);
\node[anchor=north west,font=\sffamily\fontsize{22}{24}\selectfont\bfseries,text=black!50] at (0.85,-2.35) {02};
\node[anchor=north west,text width=5.8cm,align=left,font=\sffamily\large\bfseries] at (2.0,-2.3200000000000003) {导数与梯度};
\node[anchor=north west,text width=5.9cm,align=left,font=\sffamily\small] at (2.0,-3.0700000000000003) {保留平方内层，再乘内层导数\\[.25em]偏导时固定其他变量};
\draw[black!35,line width=.4pt] (0.85,-5.050000000000001) -- (7.8999999999999995,-5.050000000000001);
\node[anchor=north west,font=\sffamily\fontsize{22}{24}\selectfont\bfseries,text=black!50] at (0.85,-7.7) {03};
\node[anchor=north west,text width=5.8cm,align=left,font=\sffamily\large\bfseries] at (2.0,-7.67) {条件概率与期望};
\node[anchor=north west,text width=5.9cm,align=left,font=\sffamily\small] at (2.0,-8.42) {条件事件决定分母\\[.25em]按相应概率加权求和};
\draw[black!35,line width=.4pt] (0.85,-10.4) -- (7.8999999999999995,-10.4);
\node[anchor=north west,font=\sffamily\fontsize{22}{24}\selectfont\bfseries,text=black!50] at (-7.35,-7.7) {04};
\node[anchor=north west,text width=5.8cm,align=left,font=\sffamily\large\bfseries] at (-6.199999999999999,-7.67) {矩阵乘法};
\node[anchor=north west,text width=5.9cm,align=left,font=\sffamily\small] at (-6.199999999999999,-8.42) {先检查内维度，再做行列点积\\[.25em]外维度决定输出形状};
\draw[black!35,line width=.4pt] (-7.35,-10.4) -- (-0.2999999999999998,-10.4);
\node[anchor=north west,font=\sffamily\fontsize{22}{24}\selectfont\bfseries,text=black!50] at (-7.35,-13.05) {05};
\node[anchor=north west,text width=5.8cm,align=left,font=\sffamily\large\bfseries] at (-6.199999999999999,-13.020000000000001) {行列式};
\node[anchor=north west,text width=5.9cm,align=left,font=\sffamily\small] at (-6.199999999999999,-13.770000000000001) {沿零项多的一行或列展开\\[.25em]方阵行列式非零时可逆};
\draw[black!35,line width=.4pt] (-7.35,-15.75) -- (-0.2999999999999998,-15.75);
\node[anchor=north west,font=\sffamily\fontsize{22}{24}\selectfont\bfseries,text=black!50] at (0.85,-13.05) {06};
\node[anchor=north west,text width=5.8cm,align=left,font=\sffamily\large\bfseries] at (2.0,-13.020000000000001) {独立验收};
\node[anchor=north west,text width=5.9cm,align=left,font=\sffamily\small] at (2.0,-13.770000000000001) {四题至少三题写全过程\\[.25em]记录第一处卡点};
\draw[black!35,line width=.4pt] (0.85,-15.75) -- (7.8999999999999995,-15.75);
\draw[->,black!55,line width=.8pt] (-.12,-3.7) -- (.55,-3.7);
\draw[->,black!55,line width=.8pt] (7.85,-5.45) -- (7.85,-7.2);
\draw[->,black!55,line width=.8pt] (.52,-9.05) -- (-.15,-9.05);
\draw[->,black!55,line width=.8pt] (-7.85,-10.8) -- (-7.85,-12.55);
\draw[->,black!55,line width=.8pt] (-.12,-14.4) -- (.55,-14.4);
\node[font=\sffamily\small,text=black!55] at (0,-18.1) {10 分钟回顾 · 25 分钟讲解 · 40 分钟练习 · 15 分钟订正};
\end{tikzpicture}
\end{center}
\clearpage

\chaptertitle{第 2 天 · 绪论与数学基础}
\textbf{2026 年 11 月 16 日｜复习模式｜K01–K02｜90 分钟}

\textbf{目标。} 能用一句话解释人工智能和图灵测试，独立完成偏导、期望、矩阵乘法和行列式四项中的至少三项。偏导标明求导变量与固定变量，乘法先标维度。

\textbf{时间。} 10 分钟回顾（求导与更新 8 分钟、图表示 2 分钟），25 分钟读绪论与四个例题，40 分钟练习，15 分钟订正。做完练习再翻答案。

\textbf{前置。} 会代入数值、做四则运算。函数写成 $f(x,y)$ 表示输入有两个分量。矩阵的 $m\times n$ 表示 $m$ 行 $n$ 列。

\section{热身}
\textbf{H1.} 新编题，沿用第 1 天错题本的再次复做题。令 $L(w)=(w+2)^2$。先令 $u=w+2$，写出 $(u^2)'$ 与 $u'(w)$ 并合成 $L'(w)$。随后取 $w_0=-1$、$\eta=0.1$，求 $L'(w_0)$ 和 $w_1=w_0-\eta L'(w_0)$，代回比较前后损失。

\textbf{H2.} 新编题。口述邻接表与邻接矩阵各存什么。若有 6 个顶点、4 条无向边，全图 DFS 分别扫描多少个邻接表条目、多少个矩阵单元？假设对每个首次访问的顶点扫完对应的表或整行。

\section{绪论速记}
人工智能研究如何让计算系统完成感知、推理、学习、规划或语言理解等任务。图灵测试的一种经典描述是让评判者只通过文字交流辨别对方是人还是机器。它考察对话表现，单凭这一表现无法判断系统内部如何推理。

机器学习关注怎样从数据中学到规律，例如从标注样本中学到分类规则。计算机视觉处理图像或视频，例如定位照片里的人。自然语言处理处理人类语言，例如翻译一段文字。领域可以重叠，例如用机器学习训练识别图片内容的视觉模型。

\clearpage
\section{单元 1 · 链式法则、偏导与极值}
若 $u=g(w)$，则 $(u^2)'=2g(w)g'(w)$。先保留整个括号，再乘括号的导数。偏导 $\partial_x f$ 表示只对 $x$ 求导，其余自变量固定。做一步梯度更新时，先求出导函数，再代入当前点，最后更新参数。

对一元函数，在内点 $f'(a)=0$ 且二阶导数存在时，$f''(a)>0$ 给出局部极小，$f''(a)<0$ 给出局部极大。$f''(a)=0$ 时要继续查一阶导数变号或比较邻域函数值。闭区间求全局极值还要把端点和不可导的候选点纳入比较。

\textbf{例题（新编）。} 设 $f(x,y)=(x-3)^2+2(y+1)^2$。求点 $(4,0)$ 的梯度，并在固定 $y=0$ 时判定关于 $x$ 的极值。

\begin{enumerate}
\item 对 $x$ 求导时固定 $y$，第二项为常数，所以 $\partial_x f=2(x-3)$。对 $y$ 求导时固定 $x$，得到 $\partial_y f=4(y+1)$。
\item 在 $(4,0)$ 代入，梯度为 $(2,4)$。两分量分别对应 $x$ 和 $y$。
\item 固定 $y=0$ 得 $(x-3)^2+2$。导数在 $x=3$ 为零，二阶导数为 2，故为局部极小。也可用平方项非负检验，全实数范围内的最小值是 2。
\end{enumerate}
\subsection{举一反三}
\CompactQuestion{A1}{固定 $y$，补齐 $\partial_x[(x+4)^2+3y^2]=2(\underline{\quad})\cdot1$。}
\CompactQuestion{A2}{令 $f(x,y)=(2x-1)^2+(y-3)^2$。求点 $(1,2)$ 的梯度。固定 $y=2$，求关于 $x$ 的极小点和极小值。指出括号的导数相较例题怎样变化。}
\clearpage
\section{单元 2 · 条件概率与期望}
当 $P(B)>0$，$P(A\mid B)=P(A\cap B)/P(B)$。条件是已知 $B$ 发生，概率的分母因此是 $P(B)$。互补事件满足 $P(A^c)=1-P(A)$，并集用 $P(A\cup B)=P(A)+P(B)-P(A\cap B)$，避免交集重复计算。

离散期望 $E[X]=\sum_x xP(X=x)$。条件期望用同样的加权求和，但权重换成条件概率。权重之和应为 1，算出的平均值应介于最小和最大结果之间。方差描述围绕均值的离散程度，$\operatorname{Var}(X)=E[X^2]-E[X]^2$。协方差 $E[(X-E[X])(Y-E[Y])]$ 描述成对变量的线性共变方向。

\textbf{例题（新编）。} 记录 10 次行动，成功 6 次且奖励都是 5，失败 4 次且奖励都是 $-1$。其中走路线 B 的有 5 次，3 次成功。把频率当成本题给定概率，求全体期望和路线 B 的条件期望。

\begin{enumerate}
\item 全体权重为 $6/10$ 与 $4/10$，分母是所有行动，所以 $E[R]=5(6/10)-1(4/10)=2.6$。
\item 在 B 的记录里成功权重为 $3/5$，其余 $2/5$ 为失败。条件期望为 $5(3/5)-1(2/5)=2.6$。
\item 两个结果相等源于本题两组成功比例恰好相同。两组权重的和都是 1，期望都在 $[-1,5]$。
\end{enumerate}
\subsection{举一反三}
\CompactQuestion{B1}{8 次行动中有 4 次走 B，其中 1 次成功。补齐 $P(\mathrm{success}\mid B)=1/\underline{\quad}$。}
\CompactQuestion{B2}{10 次行动中成功 7 次，奖励 6，失败 3 次，奖励 $-2$。其中走 B 的 4 次里只有 1 次成功。分别求全体期望和 B 的条件期望，并说明两者的权重从哪里来。}
\clearpage
\section{单元 3 · 矩阵乘法与转置}
若 $A$ 是 $m\times n$，$B$ 是 $n\times p$，则 $AB$ 有定义且为 $m\times p$。结果第 $i$ 行第 $j$ 列是 $A$ 的第 $i$ 行与 $B$ 的第 $j$ 列的点积。转置交换行和列，满足 $(AB)^T=B^TA^T$。

\textbf{例题（新编）。} $A=\begin{pmatrix}1&2&0\\0&-1&3\end{pmatrix}$，$B=\begin{pmatrix}2&1\\1&0\\0&2\end{pmatrix}$。求 $AB$。

\begin{enumerate}
\item $(2\times3)(3\times2)$ 的内维度都是 3，所以结果为 $2\times2$。
\item 第一行与两列点积得 $1(2)+2(1)+0(0)=4$，$1(1)+2(0)+0(2)=1$。第二行与两列点积得 $-1$、$6$。
\item 故 $AB=\begin{pmatrix}4&1\\-1&6\end{pmatrix}$。交换顺序虽仍可乘，但 $BA$ 为 $3\times3$。
\end{enumerate}
\subsection{举一反三}
\CompactQuestion{C1}{用例题的 $A,B$，补齐 $(AB)_{21}=0(2)+(-1)(1)+3(\underline{\quad})$，求值。}
\CompactQuestion{C2}{保留例题的 $A$，把右矩阵换成列向量 $v=(1,0,2)^T$。先标维度，再算 $Av$，并判断 $vA$ 能否定义。指出右矩阵改变后结果的列数。}
\clearpage
\section{单元 4 · 行列式与可逆}
行列式取自方阵。二阶行列式 $ad-bc$，三阶可以沿任意行或列展开。第一行的代数余子式符号依次为 $+,-,+$，每项要删去该元素所在的行和列。三角矩阵也可直接乘主对角线元素。实方阵可逆当且仅当行列式非零。

\textbf{例题（新编）。} 计算 $M=\begin{pmatrix}2&0&1\\1&3&0\\0&2&1\end{pmatrix}$ 的行列式，判断可逆性。

\begin{enumerate}
\item 沿第一行展开可跳过中间零项。第一项余子式为 $3(1)-0(2)=3$。
\item 第三项余子式为 $1(2)-3(0)=2$，符号为正。
\item 合计 $2(3)+1(2)=8$，非零，故可逆。
\end{enumerate}
\subsection{举一反三}
\CompactQuestion{D1}{用例题的 $M$，沿第一行展开，补齐第三项删去的列号以及代数余子式的符号，再求第三项的贡献。}
\CompactQuestion{D2}{令 $N=\begin{pmatrix}1&2&0\\3&5&0\\0&1&t\end{pmatrix}$。选运算量较小的一行或列展开，求 $\det N$，并给出不可逆的 $t$。指出这一题改用哪一行或列。}
\clearpage
\section{独立练习 · 约 25 分钟}
本页四题用于验收。40 分钟练习时段先用约 15 分钟完成各单元的填步和变式，再独立做以下四题。全日练习均为新编题，覆盖仓库《知识清单》K01–K02，并回顾第 1 天的 K03、K08。H1 沿用错题本的再次复做题。

\CompactQuestion{P1}{$f(x,y)=(x-2)^2+(2y+1)^2$。标出每次求导时固定的变量，求在 $(0,0)$ 的梯度。固定 $y=0$，找出关于 $x$ 的极小点及极小值。}
\CompactQuestion{P2}{某视觉系统在 12 张图上答对了 8 张。答对奖励 3，答错奖励 $-1$。阴天的 5 张图里有 2 张答对。按频率计算全体奖励期望、阴天的条件期望，并用一句话说明此任务属于哪个人工智能领域。}
\CompactQuestion{P3}{$A=\begin{pmatrix}1&-1&2\\0&3&1\end{pmatrix}$，$B=\begin{pmatrix}2&0\\1&2\\-1&1\end{pmatrix}$。标出输入和输出维度，计算 $AB$，并写出 $(AB)^T$。}
\CompactQuestion{P4}{$M=\begin{pmatrix}0&2&0\\1&3&4\\2&0&1\end{pmatrix}$。选一行展开，写出带符号的余子式，计算行列式并判断是否可逆。}
\clearpage
\section{答案与解析 · 独立作答后查看}
\subsection{热身}
\textbf{H1.} $(u^2)'=2u$（对 $u$），$u'(w)=1$，所以 $L'(w)=2(w+2)$。$L'(-1)=2$，$w_1=-1-0.1(2)=-1.2$。损失从 $1$ 降到 $0.64$。用当前点的导数决定更新量。

\textbf{H2.} 邻接表列出每个顶点的邻居，邻接矩阵逐格表示两顶点是否相邻。无向边在两端各记一次，合计 $2(4)=8$ 项。矩阵每行 6 格，共扫 $6^2=36$ 格。

\subsection{单元练习}
\textbf{A1.} 补 $x+4$。保留平方的内层，$3y^2$ 在求 $x$ 的偏导时为常数。

\textbf{A2.} $\partial_x f=2(2x-1)\cdot2=4(2x-1)$，$\partial_y f=2(y-3)$，代入是 $(4,-2)$。内层关于 $x$ 的导数改为 2。固定 $y=2$ 后，当且仅当 $x=1/2$ 时平方项为零，极小值为 1。

\textbf{B1.} 补 4，得 $1/4$。已知走 B，样本空间缩为走 B 的 4 次。

\textbf{B2.} 全体用 $7/10,3/10$，期望 $6(7/10)-2(3/10)=3.6$。B 的权重为 $1/4,3/4$，条件期望 $6(1/4)-2(3/4)=0$。两组各自归一化。

\textbf{C1.} 补 0，结果为 $-1$。取第 2 行和第 1 列。

\textbf{C2.} $(2\times3)(3\times1)$ 得 $2\times1$。两行点积为 $1$ 与 $6$，所以 $Av=(1,6)^T$。$vA$ 的内维度 1 和 2 不相等，所以未定义。结果从两列变为一列。

\textbf{D1.} 删第一行和第三列，位置 $(1,3)$ 符号为正。第三项贡献为 $1[1(2)-3(0)]=2$。

\textbf{D2.} 沿第三列展开，只有 $(3,3)$ 可能非零，符号为正。$\det N=t[1(5)-2(3)]=-t$，所以 $t=0$ 时不可逆。

\clearpage
\subsection{混合题答案}
\textbf{P1.} 固定 $y$ 得 $\partial_x f=2(x-2)$。固定 $x$ 得 $\partial_y f=4(2y+1)$，内层斜率为 2。梯度为 $(-4,4)$。固定 $y=0$ 得 $(x-2)^2+1$，在 $x=2$ 取极小值 1。

\textbf{P2.} 全体是 $3(8/12)-1(4/12)=5/3$。已知阴天，重算权重，得 $3(2/5)-1(3/5)=3/5$。输入是图片，识别属于计算机视觉。

\textbf{P3.} $(2\times3)(3\times2)$ 得 $2\times2$。四项按行乘列，依次为 $2-1-2=-1$、$0-2+2=0$、$0+3-1=2$、$0+6+1=7$。所以 $AB=\begin{pmatrix}-1&0\\2&7\end{pmatrix}$，$(AB)^T=\begin{pmatrix}-1&2\\0&7\end{pmatrix}$。

\textbf{P4.} 沿第一行，只保留第 2 列的 2。该位置符号为负，代数余子式是 $-[1(1)-4(2)]=7$，行列式为 $2(7)=14$，所以可逆。

\section{完成标准与反馈}
最后 15 分钟逐题对照，圈出第一个不同的步骤。复述四条检查规则：平方括号保留，条件发生后重定分母，乘法先对内维度，行列式展开先标符号。另闭卷回忆人工智能定义、两个领域及实例。

P1–P4 至少三题独立写全才达到计划验收。回传 H1 的推导、四题的结果和状态（独立完成／有提示完成／待复做）、总用时和第一处卡点。第 3 天会用图搜索题复测邻接表计数。

\end{document}
